\documentclass[11pt,reqno]{amsart}

%=====================================================================
%  Team A1 -- Role 1 (General Carbon Theory)
%  Flexible Computing Loads, Grid Resilience, and Carbon Emissions
%  Working results: general capacity-substitution channel and the
%  general net-carbon-reduction condition.
%=====================================================================

\usepackage{amsmath,amssymb,amsthm,mathtools}
\usepackage[margin=1in]{geometry}
\usepackage{microtype}
\usepackage{enumitem}
\usepackage{booktabs}
\usepackage{xcolor}
\usepackage[colorlinks=true,linkcolor=blue!55!black,citecolor=blue!55!black]{hyperref}

\theoremstyle{plain}
\newtheorem{theorem}{Theorem}
\newtheorem{proposition}{Proposition}
\newtheorem{lemma}{Lemma}
\newtheorem{corollary}{Corollary}
\theoremstyle{definition}
\newtheorem{condition}{Condition}
\newtheorem{assumption}{Assumption}
\theoremstyle{remark}
\newtheorem{remark}{Remark}

\newcommand{\KR}{K_R}
\newcommand{\KF}{K_F}
\newcommand{\KM}{K_M}
\newcommand{\KRs}{K_R^{*}}
\newcommand{\KFs}{K_F^{*}}
\newcommand{\KMs}{K_M^{*}}
\newcommand{\KRo}{K_R^{\circ}}
\newcommand{\KFo}{K_F^{\circ}}
\newcommand{\Dbar}{\bar D}
\newcommand{\CM}{C_M}
\newcommand{\CF}{C_F}
\newcommand{\CR}{C_R}
\newcommand{\CMbar}{\bar C_M}
\newcommand{\bp}{\beta+p}
\newcommand{\FKR}{F_{\!\KR}}
\newcommand{\E}{\mathbb{E}}
\newcommand{\PP}{\mathbb{P}}
\DeclareMathOperator*{\argmax}{arg\,max}
\newcommand{\eqdef}{\vcentcolon=}

\begin{document}

\title[General Carbon Theory for Flexible Computing Loads]
{General Carbon Theory for Flexible Computing Loads:\\
The Capacity-Substitution Channel and a Distribution-Free\\
Net-Carbon-Reduction Condition}
\author{Team A1 --- Role 1 (General Carbon Theory)}
\date{Working results, Week~1}

\begin{abstract}
We harden and generalize the carbon analysis of the flexible-computing-load
(\textsc{fcl}) model beyond the uniform-distribution benchmark. Three results.
\emph{First}, we re-derive the uniform closed forms (paper eqs.\ 18--26) from
the first-order conditions and verify them to machine precision against direct
quadrature; along the way we establish the \emph{knife-edge} structure for the
entire power-law family and reconcile two transcription issues in the current
draft. \emph{Second}, we prove the central comparative-statics claim
$\KRs>\KRo$---renewable capacity is strictly higher with the \textsc{fcl}
program---for \emph{any} residual-demand distribution satisfying a single,
economically transparent condition on the \emph{shift ratio}
$g(x)=\FKR(x)/f(x)$: the program raises renewable capacity whenever an extra
megawatt of renewables relaxes residual demand more in deep-surplus states than
in tight states. The power-law family satisfies this for every shape parameter,
which is exactly why footnote~6's extension goes through. \emph{Third}, we give
the exact net-carbon inequality, decompose it into a distribution-free
\emph{utilization} effect and a sign-controlled \emph{capacity-substitution}
effect, and state a clean sufficient condition for the carbon part of the triple
win. Every quantitative claim is checked numerically. A companion gap memo
records precisely where the fully general argument resists and what would close
it.
\end{abstract}

\maketitle

%---------------------------------------------------------------------
\section{Setup and notation}\label{sec:setup}

We work with the model of Section~3 of the paper and adopt its notation. A
regulated utility chooses renewable capacity $\KR$, dispatchable fossil backup
$\KF$, and flexible-computing-load capacity $\KM$ to maximize expected profit
over one representative period. The state variable is \emph{residual demand}
\begin{equation}\label{eq:resdem}
  \tilde x \eqdef \tilde D - \KR\tilde S \in[-\KR,\Dbar],
\end{equation}
where $\tilde D\in[0,\Dbar]$ is stochastic consumer demand and
$\tilde S\in[0,1]$ the renewable capacity factor. Let $F(x;\KR)$ be the CDF of
$\tilde x$ with density $f(x;\KR)$, extended by $F\equiv0$ for $x<-\KR$ and
$F\equiv1$ for $x>\Dbar$. More renewables shift mass left:
$\FKR(x)\eqdef\partial F/\partial\KR\ge 0$ for every $x$.

Per-MWh quantities: \textsc{fcl} revenue $R$, fossil variable cost $c$, retail
price $p$, blackout penalty $\beta$. Amortized per-MW capital costs:
$\CR,\CF,\CM$. We maintain the paper's parameter regularity.

\begin{assumption}[Regularity]\label{ass:reg}
$R>c+\CM$,\quad $R<c+\CF+\CM$,\quad $\bp>c+\CF$,\quad $\bp>R$.
\end{assumption}

Under the optimal within-period dispatch rule
$d(x)=\min\{\max\{\KF-x,0\},\KM\}$ (paper Lemma~1), plugging into expected
profit and integrating by parts gives the unified profit
\begin{equation}\label{eq:profit}
  \Pi(\KR,\KF,\KM)=p\,\E[\tilde D]
  +R\!\!\int_{\KF-\KM}^{\KF}\!\!\!F\,dx
  -c\!\!\int_{-\KM}^{\KF-\KM}\!\!\!(1-F)\,dx
  -(\bp)\!\!\int_{\KF}^{\Dbar}\!\!\!(1-F)\,dx
  -\CF\KF-\CR\KR-\CM\KM ,
\end{equation}
valid in both the large-\textsc{fcl} regime ($\KM\ge\KR$) and the
small-\textsc{fcl} regime ($\KM<\KR$) by the extended-CDF convention. Setting
$\KM=0$ gives the no-\textsc{fcl} benchmark $\Pi_0$.

\begin{lemma}[Unified first-order conditions]\label{lem:foc}
Differentiating \eqref{eq:profit} by Leibniz's rule,
\begin{align}
  \tfrac{\partial\Pi}{\partial\KM}
    &=(R-c)\,F(\KF-\KM;\KR)+c\,F(-\KM;\KR)-\CM,\label{eq:gM}\\[2pt]
  \tfrac{\partial\Pi}{\partial\KF}
    &=R\big[F(\KF)-F(\KF-\KM)\big]-c\big[1-F(\KF-\KM)\big]
      +(\bp)\big[1-F(\KF)\big]-\CF,\label{eq:gF}\\[2pt]
  \tfrac{\partial\Pi}{\partial\KR}
    &=R\!\!\int_{\KF-\KM}^{\KF}\!\!\FKR\,dx
      +c\!\!\int_{-\KM}^{\KF-\KM}\!\!\FKR\,dx
      +(\bp)\!\!\int_{\KF}^{\Dbar}\!\!\FKR\,dx-\CR.\label{eq:gR}
\end{align}
In the large-\textsc{fcl} regime $F(-\KM;\KR)=0$, and \eqref{eq:gM}--\eqref{eq:gR}
collapse to the paper's eqs.\ (12), (13), (15) respectively.\footnote{We verified
this collapse and the small-\textsc{fcl} (Appendix~G) limit numerically; the two
regimes share one formula because the extended CDF is already zero below $-\KR$.}
\end{lemma}

Two distribution-free facts pin the equilibrium ``no-blackout'' probabilities
directly from the FOCs, independent of the shape of $F$:
\begin{equation}\label{eq:cutoffs}
  F(\KFo;\KRo)=\frac{\bp-c-\CF}{\bp-c}\eqdef\varphi^{\circ},
  \qquad
  F(\KFs;\KRs)=\frac{\bp-c-\CF-\CM}{\bp-R},
  \qquad
  F(\KFs-\KMs;\KRs)=\frac{\CM}{R-c}.
\end{equation}
The first comes from the no-\textsc{fcl} FOC \eqref{eq:gF} with $\KM=0$; the
others from \eqref{eq:gM}--\eqref{eq:gF}. They give the adoption threshold
\begin{equation}\label{eq:cbar}
  \CMbar=(R-c)\,\varphi^{\circ}=\frac{(R-c)(\bp-c-\CF)}{\bp-c},
\end{equation}
the per-MW value of computing hardware equal to the margin $R-c$ times the
fraction of time spare fossil capacity is available.

%---------------------------------------------------------------------
\section{Ground truth: the uniform closed form}\label{sec:uniform}

We re-derive the uniform benchmark used to sign the carbon result, both to lock
in correctness and to fix the algebraic constants the general theory must
reproduce. Throughout this section $F(x;\KR)=\dfrac{x+\KR}{\Dbar+\KR}$ on
$[-\KR,\Dbar]$, so $f=1/Q$ and $\FKR(x)=(\Dbar-x)/Q^2$ with $Q\eqdef\Dbar+\KR$.

\begin{lemma}[Uniform optimal capacities]\label{lem:uniform}
Write $Q^{\circ}=\Dbar+\KRo$ and $Q^{*}=\Dbar+\KRs$. Then
\begin{align}
  \KFo&=\Dbar-\frac{\CF}{\bp-c}\,Q^{\circ},
  &\KRo&=\Dbar\!\left(\sqrt{\dfrac{c}{\,2\CR-\CF^{2}/(\bp-c)\,}}-1\right),
  \label{eq:noKFKR}\\[4pt]
  \KFs&=\Dbar+\frac{R-c-\CF-\CM}{\bp-R}\,Q^{*},
  &\KFs-\KMs&=\Dbar-\frac{R-\CM-c}{R-c}\,Q^{*}.
  \label{eq:fclKFKM}
\end{align}
\end{lemma}

\begin{proof}
\eqref{eq:fclKFKM} and the first identity in \eqref{eq:noKFKR} invert the cutoff
probabilities \eqref{eq:cutoffs} under the linear CDF. For $\KRo$, substitute
$\KFo$ into the renewable FOC \eqref{eq:gR} with $\KM=0$. Writing
$\gamma\eqdef\CF/(\bp-c)$ so that $\KFo=\Dbar-\gamma Q^{\circ}$ and
$\Dbar-\KFo=\gamma Q^{\circ}$, and using
$\int_a^b\FKR\,dx=\big[(b-a)(\Dbar-\tfrac{a+b}{2})\big]/Q^2$,
\[
  \underbrace{c\,\frac{\Dbar^2-\gamma^2 Q^{\circ2}}{2Q^{\circ2}}}_{\text{from }(0,\KFo)}
  +\underbrace{\frac{\bp}{2}\gamma^2}_{\text{from }(\KFo,\Dbar)}=\CR
  \;\Longrightarrow\;
  \frac{c\,\Dbar^2}{Q^{\circ 2}}=2\CR-\frac{\CF^2}{\bp-c},
\]
because $\gamma^2(\bp-c)=\CF^2/(\bp-c)$. Solving for $Q^{\circ}$ and subtracting
$\Dbar$ gives \eqref{eq:noKFKR}. \emph{(Verified to machine precision against the
quadrature solver.)}
\end{proof}

\paragraph{Carbon (paper eqs.\ 25--26).}
Expected fossil output without the program is
$\E[\text{fossil}^{\circ}]=\KFo-\int_0^{\KFo}F\,dx
   =\KFo\big(\Dbar-\tfrac{\KFo}{2}\big)/Q^{\circ}$,
and with the program (large-\textsc{fcl} regime), accounting for the flat stretch
$(\KMs-\KRs)$ on which the extended CDF equals zero,
\begin{equation}\label{eq:fosfcl}
  \E[\text{fossil}^{*}]
  =(\KMs-\KRs)+\int_{-\KRs}^{\KFs-\KMs}\!(1-F)\,dx
  =\Dbar+\frac{R-c-\CF-\CM}{\bp-R}\,Q^{*}-\frac12\Big(\frac{\CM}{R-c}\Big)^{2}Q^{*}.
\end{equation}
Carbon emissions equal $c$ times these quantities. Net carbon falls with the
program iff $\E[\text{fossil}^{*}]<\E[\text{fossil}^{\circ}]$.

\begin{remark}[Two corrections to the current draft]\label{rem:corr}
\emph{(i)} The bracketed expressions printed in eqs.\ (25)--(26) are
\emph{expected fossil output} ($=\text{carbon}/c$), not carbon itself; the
constant $c$ is dropped because only the sign of (25)$-$(26) matters. Restoring
$c$, our \eqref{eq:fosfcl} reproduces the printed eq.\ (26) to $6\times10^{-10}$
and eq.\ (25) to machine precision.
\emph{(ii)} The naive ``small-\textsc{fcl}'' fossil formula
$\KF(\Dbar+\KM-\KF/2)/Q$ is invalid here: in the large-\textsc{fcl} regime the
integration range $[-\KM,\KF-\KM]$ protrudes below $-\KR$, where $1-F=1$
exactly. The flat term $(\KMs-\KRs)$ in \eqref{eq:fosfcl} is the correction.
\end{remark}

\begin{proposition}[Knife-edge for the power-law family]\label{prop:knife}
Let $F(x;\KR)=\big((x+\KR)/(\Dbar+\KR)\big)^{\alpha}$ on $[-\KR,\Dbar]$ for any
$\alpha>0$ (the uniform is $\alpha=1$). Along the locus where $\KF,\KM$ are
re-optimized at every $\KR$, the marginal value of renewable capacity is
\emph{constant in $\KR$}. In the interior regime it equals
\begin{equation}\label{eq:mv}
  \mathrm{MV}=\tfrac12\!\left[\,c
  +\frac{(R-\CM-c)^2}{R-c}
  +\frac{(R-c-\CF-\CM)^2}{\bp-R}\right],
\end{equation}
so $\KRs=\infty$ if $\mathrm{MV}>\CR$, $\KRs=0$ if $\mathrm{MV}<\CR$, and any
$\KR$ is optimal at equality. Equivalently $\KRs=\infty$ iff
\begin{equation}\label{eq:eq24}
  c+\frac{(R-\CM-c)^2}{R-c}+\frac{(R-c-\CF-\CM)^2}{\bp-R}>2\CR ,
\end{equation}
which is \emph{exactly} the bracket printed in paper eq.\ (24).
\end{proposition}

\begin{proof}
For the power law one computes $\FKR(x)=f(x)\,(\Dbar-x)/Q$, hence the shift ratio
$g(x)=\FKR/f=(\Dbar-x)/Q$ is the \emph{same} affine function for every $\alpha$.
Substituting the linear loci \eqref{eq:fclKFKM} into \eqref{eq:gR} and evaluating
$\int_a^b\FKR\,dx$ over the three regions (whose endpoints are all of the form
$\Dbar-(\,\cdot\,)Q$), every $Q$ cancels, leaving \eqref{eq:mv}; the marginal
return is constant in $\KR$, so the objective is linear along the locus and the
optimum is a corner. Multiplying \eqref{eq:mv} by $2$ and subtracting $2\CR$
gives precisely the paper's eq.\ (24) bracket, confirming
$\text{(eq.\ 24)}=2\,(\mathrm{MV}-\CR)=2\,\partial\Pi/\partial\KR$ along the
locus. \emph{(MV constant in $\KR$ verified at $\KR\in\{1,5,9\}$ for
$\alpha\in\{0.5,1,2\}$.)}
\end{proof}

\begin{remark}[On eq.\ (24)]
Proposition~\ref{prop:knife} shows eq.\ (24) is \emph{correct as a sign
condition}: the printed bracket equals twice the profit gradient
$\partial\Pi/\partial\KR$, and only its sign is used (the $\infty/0$
trichotomy). The factor of two is a grouping artifact of the closed-form
substitution, not an error---worth a one-line clarifying footnote in the draft.
\end{remark}

%---------------------------------------------------------------------
\section{The capacity-substitution channel, in general}\label{sec:capacity}

The carbon result rests on a single comparative-statics claim that the paper
asserts but does not prove outside the uniform case: \emph{renewable capacity is
higher with the program}, $\KRs>\KRo$. We now prove it for any distribution
satisfying one transparent condition. The obstruction, and why the naive
lattice argument fails, is recorded in Section~\ref{sec:gap}.

\subsection{The marginal-value cross-partials}
From Lemma~\ref{lem:foc}, the second derivatives of $\Pi$ that govern how the
optimum moves when computing becomes more viable are
\begin{align}
  \frac{\partial^2\Pi}{\partial\KR\,\partial\KM}
    &=(R-c)\,\FKR(\KF-\KM)+c\,\FKR(-\KM)\;\ge 0,\label{eq:hRM}\\
  \frac{\partial^2\Pi}{\partial\KF\,\partial\KM}
    &=(R-c)\,f(\KF-\KM)\;\ge 0,\label{eq:hFM}\\
  \frac{\partial^2\Pi}{\partial\KR\,\partial\KF}
    &=(R-\bp)\,\FKR(\KF)-(R-c)\,\FKR(\KF-\KM)\;\le 0.\label{eq:hRF}
\end{align}
So computing raises the marginal value of \emph{both} capacities
(\eqref{eq:hRM}--\eqref{eq:hFM} are nonnegative because $R>c$ and $\FKR\ge0$),
yet renewables and fossil are mutual substitutes in the objective
(\eqref{eq:hRF} is nonpositive because $\bp>R$). The sign conflict in
\eqref{eq:hRF} is exactly why a one-line Topkis argument does not deliver
$\KRs>\KRo$; see Section~\ref{sec:gap}.

\subsection{The shift-ratio condition}
Define the \emph{shift ratio}
\begin{equation}\label{eq:shift}
  g(x;\KR)\eqdef\frac{\FKR(x;\KR)}{f(x;\KR)}
  =-\left.\frac{\partial x}{\partial\KR}\right|_{F\,\text{fixed}} ,
\end{equation}
the megawatts by which one extra MW of renewables pulls down the residual-demand
level holding its quantile fixed---the strength of the renewable
\emph{leftward shift} at state $x$.

\begin{condition}[Shift ratio]\label{cond:cs}
$g(x;\KR)$ is non-increasing in $x$ on $(\KFs-\KMs,\,\KFs)$ at each point of the
equilibrium path $\CM\in[\CM,\CMbar]$.
\end{condition}

Condition~\ref{cond:cs} says the renewable shift bites harder in deep-surplus
states than in tight states. It is mild and interpretable, and it holds for the
entire power-law family, where $g(x)=(\Dbar-x)/Q$ is affine and strictly
decreasing for every $\alpha$.

\begin{theorem}[Capacity-substitution channel]\label{thm:capacity}
Suppose the optimum is interior with negative-definite Hessian along the path
$\CM\in[\CM,\CMbar]$, and Condition~\ref{cond:cs} holds there. Then
\begin{equation}\label{eq:dKRdCM}
  \frac{d\KRs}{d\CM}
  =\frac{(\bp-R)(R-c)\big[\FKR(\KFs)\,f(\KFs-\KMs)-\FKR(\KFs-\KMs)\,f(\KFs)\big]}
        {\det \mathcal H}\;\le 0,
\end{equation}
where $\mathcal H$ is the Hessian of $\Pi$ (so $\det\mathcal H<0$). Consequently,
integrating from the adoption boundary $\CMbar$---where $\KMs=0$ and the system
collapses to the no-\textsc{fcl} optimum, $\KRs(\CMbar)=\KRo$---we obtain
\begin{equation}
  \KRs(\CM)\;\ge\;\KRo\qquad\text{for all }\CM\le\CMbar,
\end{equation}
with strict inequality whenever $\CM<\CMbar$ and $g$ is strictly decreasing.
Renewable capacity is strictly higher with the \textsc{fcl} program.
\end{theorem}

\begin{proof}
Only $\partial\Pi/\partial\KM$ in \eqref{eq:gM} depends on $\CM$, with
$\partial^2\Pi/\partial\KM\partial\CM=-1$. Implicit differentiation of the FOC
system $\nabla\Pi=0$ gives
$d\KRs/d\CM=(\mathcal H^{-1})_{\KR\KM}
   =\operatorname{cof}_{\KM\KR}(\mathcal H)/\det\mathcal H$.
In the large-\textsc{fcl} regime $\FKR(-\KM)=f(-\KM)=0$, so by
\eqref{eq:hRM}--\eqref{eq:hRF} and $\partial^2\Pi/\partial\KM^2=-(R-c)f(\KF-\KM)$,
the relevant cofactor simplifies (after cancellation) to
\[
  \operatorname{cof}_{\KM\KR}(\mathcal H)
  =(\bp-R)(R-c)\big[\FKR(\KFs)f(\KFs-\KMs)-\FKR(\KFs-\KMs)f(\KFs)\big].
\]
Dividing by $\det\mathcal H<0$ gives \eqref{eq:dKRdCM}. The bracket equals
$f(\KFs)f(\KFs-\KMs)\big[g(\KFs)-g(\KFs-\KMs)\big]\le 0$ under
Condition~\ref{cond:cs} (since $\KFs-\KMs<\KFs$), and $(\bp-R)(R-c)>0$, so the
numerator is $\le0$ and $d\KRs/d\CM\le0$. At $\CM=\CMbar$ the marginal profit of
the first \textsc{fcl} unit \eqref{eq:gM} vanishes
($(R-c)\varphi^{\circ}-\CMbar=0$), so $\KMs=0$ and the remaining FOCs are the
no-\textsc{fcl} conditions, giving $\KRs(\CMbar)=\KRo$. Integrating the
nonpositive derivative from $\CMbar$ down to $\CM$ yields $\KRs(\CM)\ge\KRo$.
\end{proof}

\begin{corollary}[Power-law / footnote~6]\label{cor:power}
For every power-law residual demand ($\alpha>0$), $g(x)=(\Dbar-x)/Q$ is strictly
decreasing, so Condition~\ref{cond:cs} holds and $\KRs>\KRo$ whenever
$\CM<\CMbar$. This is precisely why footnote~6's closed-form extension carries
through for the whole family rather than the uniform case alone.
\end{corollary}

\begin{remark}[Reading of Condition~\ref{cond:cs}]
Writing $\tilde x=\tilde D-\KR\tilde S$, the shift ratio is the
density-weighted conditional capacity factor,
$g(x)=\E\!\big[\tilde S \mid \text{quantile at }x\big]$ in the natural sense.
Condition~\ref{cond:cs} fails only if high-demand states are systematically
\emph{windier/sunnier} than surplus states---an empirically unusual positive
demand--supply correlation. Under the more realistic negative or zero
correlation, the condition holds and the channel operates.
\end{remark}

%---------------------------------------------------------------------
\section{The net carbon condition, in general}\label{sec:carbon}

Carbon emissions equal $c$ times expected fossil output, so the program lowers
carbon iff $\E[\text{fossil}^{*}]<\E[\text{fossil}^{\circ}]$. Decompose
\eqref{eq:profit}'s fossil terms exactly. Let
$q^{\circ}\eqdef F(0;\KRo)$ be the no-\textsc{fcl} surplus probability and define
the conditional means
\[
  \mu^{*}_{\mathrm{low}}\eqdef
  \frac{1}{F(\KFs-\KMs;\KRs)}\!\int_{-\KRs}^{\KFs-\KMs}\!\!x\,f(x;\KRs)\,dx,
  \qquad
  \mu^{\circ}_{\mathrm{mid}}\eqdef
  \frac{1}{\varphi^{\circ}-q^{\circ}}\!\int_{0}^{\KFo}\!\!x\,f(x;\KRo)\,dx .
\]

\begin{proposition}[Exact net-carbon inequality]\label{prop:carbon}
For any $F$ satisfying Assumption~\ref{ass:reg},
\begin{align}
  \E[\text{fossil}^{*}] &=\KFs-\tfrac{\CM}{R-c}\big(\KFs-\KMs-\mu^{*}_{\mathrm{low}}\big),
  \label{eq:fos1}\\
  \E[\text{fossil}^{\circ}] &=\KFo(1-\varphi^{\circ})+(\varphi^{\circ}-q^{\circ})\,\mu^{\circ}_{\mathrm{mid}},
  \label{eq:fos2}
\end{align}
and carbon falls with the program iff
\begin{equation}\label{eq:C}
  \underbrace{\frac{\CM}{R-c}\big(\KFs-\KMs-\mu^{*}_{\mathrm{low}}\big)}_{\text{utilization effect }(+)}
  \;>\;
  \underbrace{\KFs-\KFo(1-\varphi^{\circ})-(\varphi^{\circ}-q^{\circ})\mu^{\circ}_{\mathrm{mid}}}_{\text{net of capacity effects}}.
  \tag{C}
\end{equation}
\end{proposition}

\begin{proof}
Integrate each fossil integral by parts using
$\int_a^b(1-F)\,dx=[x(1-F)]_a^b+\int_a^b x f\,dx$; for \eqref{eq:fos1} take
$a=-\KMs,\ b=\KFs-\KMs$ with $F(a;\KRs)=0$ (large-\textsc{fcl}) and
$F(b;\KRs)=\CM/(R-c)$; for \eqref{eq:fos2} take $a=0,\ b=\KFo$ with
$F(\KFo;\KRo)=\varphi^{\circ}$. Subtract.
\end{proof}

Two forces. The \textbf{utilization effect} (the left side of (C)) is
\emph{always positive and distribution-free}: at fixed capacities the program
strictly raises fossil output in every state below $\KF$ (Lemma~1 dispatch).
The \textbf{capacity-substitution effect} is carried by the right side through
$\KRs>\KRo$ (Theorem~\ref{thm:capacity}), which lowers $\E[\text{fossil}^{*}]$ by
shifting the residual-demand distribution left. The net sign is which one wins.

\begin{corollary}[Sufficient distributional conditions]\label{cor:suff}
The carbon part of the triple win holds for all $\CM<\CMbar$ under any $F$ with
\begin{enumerate}[label=(\alph*),leftmargin=2.4em]
  \item $f(\cdot;\KRs)$ non-increasing on $(-\KRs,\KFs-\KMs)$ \;(equiv.\ $F(\cdot;\KRs)$ concave there), and
  \item $f(\cdot;\KRo)$ non-decreasing on $(0,\KFo)$ \;(equiv.\ $F(\cdot;\KRo)$ convex there).
\end{enumerate}
Under (a), $\mu^{*}_{\mathrm{low}}\le(\KFs-\KMs-\KRs)/2$, so \textsc{fcl} fossil
output is no larger than under the uniform; under (b),
$\mu^{\circ}_{\mathrm{mid}}\ge\KFo/2$, so no-\textsc{fcl} output is no smaller.
Together (C) holds at least as strongly as under the uniform, and the threshold
$\CMbar$ is preserved. The uniform satisfies (a)--(b) with equality.
\end{corollary}

\begin{theorem}[Clean sufficient condition for the carbon triple win]\label{thm:carbon}
Suppose Condition~\ref{cond:cs} holds (so $\KRs>\KRo$) and the residual-demand
density is unimodal with mode in the positive-demand region. Then there is a
threshold $\CR^{\dagger}>0$ such that, for all $\CR\le\CR^{\dagger}$, the program
strictly lowers carbon for every $\CM<\CMbar$. As $\CR\downarrow$, the
capacity-substitution effect grows without bound---renewable investment expands,
residual demand moves into deep surplus, and fossil output (hence carbon)
$\to 0$---while the utilization effect stays bounded by
$\tfrac{\CM}{R-c}(\KFs-\KMs)$. The two cross at $\CR^{\dagger}$.
\end{theorem}

\noindent Under the uniform/power-law family this is the knife-edge of
Proposition~\ref{prop:knife}: when \eqref{eq:eq24} holds, $\KRs=\infty$, fossil
output vanishes, and carbon strictly falls---recovering paper eqs.\ (25)--(26).

%---------------------------------------------------------------------
\section{Numerical validation}\label{sec:num}

All claims are checked by an independent solver that quadratures the unified
profit \eqref{eq:profit} and solves \eqref{eq:gM}--\eqref{eq:gR} by Newton
methods (no closed forms assumed). Highlights:
\begin{itemize}[leftmargin=1.6em]
  \item \textbf{Uniform closed forms.} Carbon eqs.\ (25)/(26) match quadrature to
        $\le 6\times10^{-10}$; capacities \eqref{eq:noKFKR}--\eqref{eq:fclKFKM}
        satisfy the FOCs to $\le 10^{-15}$.
  \item \textbf{Knife-edge.} $\mathrm{MV}(\KR)$ constant across $\KR\in\{1,5,9\}$
        for $\alpha\in\{0.5,1,2\}$; eq.\ \eqref{eq:eq24} reproduces paper eq.\ (24).
  \item \textbf{Capacity channel.} On strictly-concave Beta scale-mixtures
        $\tilde x=\tilde D-\KR\tilde S$ with interior optima, across symmetric,
        decreasing-density, and right-skewed consumer demand
        ($D\sim\Dbar\cdot\mathrm{Beta}(2,2),(1,3),(2,5),(2,8)$), every case has
        $g$ strictly decreasing, $\KRs>\KRo$, and the sign of $d\KRs/d\CM$
        matching the shift-ratio bracket of \eqref{eq:dKRdCM}. Representative
        cell ($\mathrm{Beta}(1,3)$): $\KRs=3.274>\KRo=0.831$,
        $g(\KFs-\KMs)=0.457>g(\KFs)=0.414$, $d\KRs/d\CM=-2.99<0$ with bracket
        $-3.1\times10^{-4}<0$. Strongly left-skewed demand pushes the optimum to
        the full-coverage boundary ($\KFs\!\to\!\Dbar$), i.e.\ footnote~6's
        ``harder'' regime, where the interior-Hessian hypothesis of
        Theorem~\ref{thm:capacity} no longer applies.
  \item \textbf{Sharpness.} The bracket in \eqref{eq:dKRdCM} changes sign with
        $g(\KFs)-g(\KFs-\KMs)$, so $g$ \emph{increasing} reverses the channel
        ($\KRs<\KRo$); the shift-ratio condition is therefore tight, not merely
        sufficient.
\end{itemize}

%---------------------------------------------------------------------
\section{Where the general argument resists (gap memo)}\label{sec:gap}

A self-contained one-pager accompanies this note (\texttt{gap\_memo.md}). In
brief:
\begin{enumerate}[leftmargin=1.7em]
  \item \textbf{Submodularity obstruction.} By \eqref{eq:hRF}, $\Pi$ is
        \emph{submodular} in $(\KR,\KF)$ while having increasing differences in
        each of $(\KR;\KM)$ and $(\KF;\KM)$. A single change of order cannot fix
        both signs, so Topkis/Milgrom--Shannon does not directly yield
        $\KRs>\KRo$. Theorem~\ref{thm:capacity} sidesteps this with a
        homotopy in $\CM$ plus the local cofactor sign---valid under
        interiority and a negative-definite Hessian along the path.
  \item \textbf{Endogenous-capacity feedback in (C).} Corollary~\ref{cor:suff}
        compares conditional means \emph{holding the FOC-pinned cutoff
        probabilities fixed}. A fully general statement must also track how
        $\KRs,\KFs$ themselves move with the shape of $F$; that second-order
        feedback is not yet signed in closed form.
  \item \textbf{$\CR^{\dagger}$ not explicit.} Theorem~\ref{thm:carbon}
        establishes existence of the renewable-cost threshold but not a closed
        form outside the power-law family.
  \item \textbf{Boundary regimes.} The proof assumes the large-\textsc{fcl}
        interior regime along the whole path; reconciling the small-\textsc{fcl}
        (Appendix~G) and full-coverage ($R\ge c+\CF+\CM$) boundaries into one
        statement is left to Role~4.
\end{enumerate}

\end{document}
